% ─────────────────────────────────────────────────────────────────────────────
\chapter{Repository and Reproducibility Artifacts}
\label{app:repository-reproducibility-artifacts}
% ─────────────────────────────────────────────────────────────────────────────

This appendix documents the technical artifacts that support the reproducibility,
verifiability, and traceability of the experimental work described in the
thesis. It does not replace the methodology presented in
Chapter~\ref{chap:methodology}, nor does it anticipate the analysis of the
results. Rather, it provides an operational reference for connecting the
academic description of the pipeline to the scripts, manifests, hashes, and
outputs maintained in the repository associated with the project. This approach
is consistent with the general principles of \gls{df}, in which the
preservation of provenance, artifact integrity, and documentation of operations
are essential requirements for making a digital analysis process
verifiable~\cite{casey2011digital,garfinkel2010digital,nist2006guide,iso27037}.

The complete source code is not reported in full within the document, since
duplicating it in the thesis would make the text redundant and less
maintainable. The repository instead represents the appropriate location for
consulting the scripts, manifests, and operational documentation. The appendix
therefore describes the logical structure of the project, the official artifacts
of the pipeline, the main execution sequence, the traceability mechanisms, and
the boundaries within which reproducibility must be interpreted. In particular,
reproducibility concerns the controlled reconstruction of the experimental
artifacts starting from the official manifests and scripts, whereas it does not
necessarily imply the public redistribution of all raw images, which may be
limited by legal, ethical, or source-dependent constraints.


% ─────────────────────────────────────────────────────────────────────────────
\section{Repository Structure}
\label{app:repository-structure}
% ─────────────────────────────────────────────────────────────────────────────

The repository is organized according to a functional separation among raw data,
prepared data, the final dataset, experimental splits, perturbations, proxy
models, evaluation, explainability, results, and thesis documentation. This
organization makes it possible to distinguish the phases of acquisition,
preparation, manual selection, split generation, attack production, automated
evaluation, and preparation of the bundle intended for selected black-box
software systems.

\begin{lstlisting}[caption={Simplified structure of the experimental repository.}, label={lst:repository-structure}]
msc-thesis-ai-robustness-in-digital-forensics/
  datasets/
    raw/
    prepared/
    final/
    splits/
    forensic_evaluation_bundle/
    scripts/
      acquisition/
      prepared/
      final/
      splits/
      attacks/
      bundle/
      utils/
  attacks/
    adversarial/
    anti_forensic/
    manifests/
  models/
    scripts/
    checkpoints/
    reports/
  evaluation/
    scripts/
    proxy_models/
    forensic_tools/
  forensic_tools/
  explainability/
    scripts/
    outputs/
    manifests/
  results/
    scripts/
    metrics/
    tables/
    plots/
    figures/
    reports/
  docs/
    LatexThesis/
  .github/
\end{lstlisting}

The \codepath{datasets/} directory contains the data pipeline, from the initial
collection phase to the forensic evaluation bundle. The \codepath{attacks/}
directory collects the outputs of adversarial perturbations and anti-forensic
transformations, whereas \codepath{models/} contains the scripts and checkpoints
of the proxy models. The \codepath{evaluation/} directory collects the scripts
and outputs of automated evaluation and normalization of the exports produced by
selected black-box software systems. The \codepath{forensic_tools/} directory
contains the normalized black-box evaluation perimeter and related exported
artifacts, whereas the \codepath{results/scripts/} directory contains the
reporting and audit scripts used to generate thesis-ready figures, tables, and
sensitivity checks. The \codepath{explainability/} directory is reserved for the
generation of \gls{xai}-based case studies, whereas
\codepath{docs/LatexThesis/} contains the official LaTeX version of the thesis.


% ─────────────────────────────────────────────────────────────────────────────
\section{Official Experimental Artifacts}
\label{app:official-experimental-artifacts}
% ─────────────────────────────────────────────────────────────────────────────

Table~\ref{tab:official-experimental-artifacts} summarizes the main artifacts
of the pipeline. These files are not merely intermediate outputs, but constitute
the anchoring points between the methodology described in the body of the thesis
and its experimental implementation. In particular, the manifests preserve the
connection among image identifier, source, final label, hash, split, possible
perturbation, and path of the generated file.

\begin{longtable}{@{}%
>{\raggedright\arraybackslash}p{0.43\textwidth}
>{\raggedright\arraybackslash}p{0.53\textwidth}@{}}
\caption{Main official artifacts of the experimental pipeline.}
\label{tab:official-experimental-artifacts} \\
\toprule
\textbf{Artifact} & \textbf{Role in the pipeline} \\
\midrule
\endfirsthead

\toprule
\textbf{Artifact} & \textbf{Role in the pipeline} \\
\midrule
\endhead

\midrule
\multicolumn{2}{r}{\footnotesize Continued on the next page} \\
\endfoot

\bottomrule
\endlastfoot

\footnotesize\codepath{datasets/prepared/final_pool/metadata.csv}
&
Technical metadata of the prepared pool, generated after image validation, hash
computation, and removal of exact duplicates. \\

\footnotesize\codepath{datasets/prepared/manifests/review_manifest_full.csv}
&
Complete manifest used as the basis for the \textit{human-in-the-loop} manual
review protocol. \\

\footnotesize\codepath{datasets/final/manifests/manual_selection_protocol_db.csv}
&
Operational database of the manual selection process, including review status
and semantic assignments. \\

\footnotesize\codepath{datasets/final/manifests/manual_selection_final_1500.csv}
&
Official manifest of the frozen dataset, composed of 1\,500 images balanced
across the \texttt{weapon}, \texttt{non\_weapon}, and \gls{ood} evaluation
categories. \\

\footnotesize\codepath{datasets/final/manifests/manual_selection_adversarial_subset.csv}
&
Official manifest of the binary subset, composed of the \texttt{weapon} and
\texttt{non\_weapon} classes, used for clean splits, adversarial attacks, and
anti-forensic transformations. \\

\footnotesize\codepath{datasets/final/manifests/manual_selection_removed.csv}
&
Manifest of the samples excluded during the final selection, useful for auditing
and checking review decisions. \\

\footnotesize\codepath{datasets/splits/manifests/clean_folds_manifest.csv}
&
Manifest of the five clean folds, with traceability by class, source, path,
hash, and sample type. \\

\footnotesize\codepath{datasets/splits/manifests/ood_eval_manifest.csv}
&
Manifest of the \gls{ood} set, kept separate from the binary subset and used for
out-of-distribution evaluation. \\

\footnotesize\codepath{attacks/manifests/}
&
Directory of the manifests produced to connect each perturbed image to the clean
sample, fold, true class, and type of attack or transformation. \\

\footnotesize\codepath{datasets/forensic_evaluation_bundle/metadata/bundle_manifest.csv}
&
Global manifest of the bundle intended for evaluation with selected black-box
software systems. \\

\footnotesize\codepath{datasets/forensic_evaluation_bundle/metadata/bundle_hashes_sha256.csv}
&
Control file containing the \gls{sha256} hashes of the artifacts included in the
forensic bundle. \\

\footnotesize\codepath{datasets/forensic_evaluation_bundle/metadata/bundle_summary.json}
&
Structured summary of the bundle, including counts, consistency checks, and
operational instructions for import into selected black-box software systems. \\

\footnotesize\codepath{datasets/forensic_evaluation_bundle/metadata/embedded_metadata_audit.csv}
&
Audit file generated during bundle construction to document residual embedded
metadata in the blind input view. \\

\footnotesize\codepath{results/figures/chapter_5/chapter5_figures_manifest.csv}
&
Manifest of the thesis-ready figures generated from the consolidated metric
files by the reporting layer. \\

\footnotesize\codepath{results/figures/chapter_5/chapter5_figures_summary.json}
&
Structured summary of the experimental reporting assets generated for
Chapter~\ref{chap:results}. \\

\footnotesize\codepath{results/figures/chapter_5/embedded_metadata_sensitivity_summary.json}
&
Structured summary of the embedded-metadata sensitivity check used to support
the audit discussion in Chapter~\ref{chap:results}. \\

\end{longtable}


% ─────────────────────────────────────────────────────────────────────────────
\section{Dataset Traceability and Human-in-the-Loop Selection}
\label{app:dataset-traceability-human-in-the-loop}
% ─────────────────────────────────────────────────────────────────────────────

The data pipeline is based on a combination of automatic checks and manual
review. Images acquired from the original sources are technically validated,
deduplicated through hashes, and inserted into a prepared pool. Only afterward
are they included in the review manifest, in which the final classification is
assigned through a \textit{human-in-the-loop} protocol. This methodological
choice is consistent with the forensic purpose of the work: the dataset is not
treated as a purely automatic product, but as a verifiable, frozen, and
documented experimental artifact.

The final dataset is represented by the
\codepath{manual_selection_final_1500.csv} manifest and contains 1\,500 images
distributed in a balanced manner across three evaluation categories:
\texttt{weapon}, \texttt{non\_weapon}, and \gls{ood}. The \texttt{weapon} class
includes only realistic individual weapons that are clearly visible; the
\texttt{non\_weapon} class includes realistic negative samples consistent with
the task; and the \gls{ood} category collects borderline or
out-of-distribution cases, such as knives, swords, replicas, airsoft items,
synthetic renderings, video game images, heavy military vehicles, explosions,
war scenes, and degraded or anomalous images. The official adversarial subset,
stored in \codepath{manual_selection_adversarial_subset.csv}, excludes
\gls{ood} examples and retains only the two operational classes
\texttt{weapon} and \texttt{non\_weapon}, for a total of 1\,000 images.

Stable identifiers, \gls{sha256} hashes, relative paths, and intermediate
manifests make it possible to reconstruct the derivation of each sample from
the prepared pool to the splits, perturbed versions, and evaluation outputs.
This connection is particularly relevant in the black-box evaluation phase,
since selected software systems may rename files, modify the structure of
exports, or produce reports in heterogeneous formats. The use of hashes
therefore makes it possible to maintain a criterion that is independent of the
file name and export path.


% ─────────────────────────────────────────────────────────────────────────────
\section{Official Script Sequence}
\label{app:official-script-sequence}
% ─────────────────────────────────────────────────────────────────────────────

The operational pipeline uses numbered scripts as official entry points. This
numbering documents the logical order of the operations and reduces ambiguity in
the reconstruction of the workflow. Listing~\ref{lst:official-script-sequence}
reports the main sequence aligned with the current state of the repository.

\begin{lstlisting}[caption={Official sequence of the main pipeline scripts.}, label={lst:official-script-sequence}]
# Raw acquisition and source preparation
python datasets/scripts/acquisition/00_download_raw_datasets_bundle.py
python datasets/scripts/acquisition/01_download_kaggle.py
python datasets/scripts/acquisition/02_download_github.py
python datasets/scripts/acquisition/03_build_subset_deepfirearm.py
python datasets/scripts/acquisition/04_scrape_google.py
python datasets/scripts/acquisition/05_scrape_telegram.py
python datasets/scripts/acquisition/06_scrape_youtube.py
python datasets/scripts/acquisition/07_scrape_deepweb.py

# Prepared dataset and human-in-the-loop selection
python datasets/scripts/prepared/08_build_prepared_dataset.py
python datasets/scripts/prepared/09_generate_review_manifest_full.py
python datasets/scripts/final/10_manual_selection_protocol_reviewer.py

# Clean/OOD split generation
python datasets/scripts/splits/11_generate_clean_and_ood_splits.py

# Proxy model training and perturbation generation
python models/scripts/12_train_proxy_models.py
python datasets/scripts/attacks/13_generate_anti_forensic_attacks.py
python datasets/scripts/attacks/14_generate_adversarial_attacks.py

# Proxy model evaluation and forensic bundle construction
python evaluation/scripts/15_evaluate_proxy_models.py
python datasets/scripts/bundle/16_build_forensic_evaluation_bundle.py

# Explainability case-study preparation
python explainability/scripts/17_generate_integrated_gradients_case_studies.py
python explainability/scripts/18_xai_interactive_launcher.py

# Black-box software export normalization
python evaluation/scripts/19_normalize_forensic_ai_tool_predictions.py

# Experimental reporting assets and metadata sensitivity
python results/scripts/20_generate_experimental_reporting_assets.py
python results/scripts/21_generate_embedded_metadata_sensitivity_check.py
\end{lstlisting}

The script
\codepath{evaluation/scripts/19_normalize_forensic_ai_tool_predictions.py}
constitutes the connection point between the exports produced by selected
black-box software systems and the final comparative metrics. In the final
protocol, it exclusively supports outputs attributable to
\texttt{magnet\_axiom}, \texttt{excire\_foto\_2025},
\texttt{cellebrite\_inseyets}, and \texttt{griffeye}. Its role is to normalize
the exports, map predictions back to the records of the forensic evaluation
bundle, and produce results that are comparable with the proxy level through
mapping based, in preferential order, on \texttt{bundle\_id}, \gls{sha256}
hash, \gls{md5} hash, file name, and manual audit.

The final two scripts belong to the reporting and audit layer of the pipeline.
The script
\codepath{results/scripts/20_generate_experimental_reporting_assets.py} reads
the consolidated metric files and produces thesis-ready figures, tables,
manifest files, and summaries for Chapter~\ref{chap:results}. The script
\codepath{results/scripts/21_generate_embedded_metadata_sensitivity_check.py}
connects the embedded-metadata audit generated during bundle construction with
the normalized black-box predictions, producing the sensitivity tables and
summary used to document the residual metadata limitation.


% ─────────────────────────────────────────────────────────────────────────────
\section{Models, Attacks, and Evaluation Outputs}
\label{app:models-attacks-evaluation-outputs}
% ─────────────────────────────────────────────────────────────────────────────

The consolidated experimental protocol uses three proxy models:
\texttt{efficientnet\_b0}, \texttt{resnet18}, and \gls{clip}. These models are
used as transparent and inspectable baselines to observe, in a controlled
manner, the behavior of classifiers with respect to clean, perturbed, and
out-of-distribution inputs. Any preliminary components or exploratory baselines
not included in the final evaluation are not considered part of the consolidated
experimental protocol. The use of proxy models places the analysis within the
broader context of robustness in \gls{ml} and \gls{ai}
systems~\cite{bishop2006pattern,biggio2018wild,radford2021learning}.

The adversarial attacks generated in the pipeline are \texttt{fgsm},
\texttt{superdeepfool}, \texttt{sigma\_zero}, \texttt{one\_pixel}, and
\texttt{color\_shift}. The anti-forensic transformations are
\texttt{jpeg\_recompression}, \texttt{resample\_resize},
\texttt{gaussian\_blur}, \texttt{histogram\_modification}, and
\texttt{contrast\_stretching}. This distinction is methodologically important:
adversarial attacks are used to evaluate the decision vulnerability of the
models, whereas anti-forensic transformations simulate operationally plausible
manipulations of the image content or file that may alter the visibility,
quality, or technical traces of the artifact~\cite{stamm2010anti,nowroozi2021survey,croce2020robustbench}.

The outputs of the proxy model evaluation are collected in the
\codepath{evaluation/proxy_models/} and \codepath{results/metrics/}
directories. These outputs include predictions, aggregate metrics, confusion
matrices, and comparisons among clean, adversarial, anti-forensic, and
\gls{ood} conditions. Their function is not only quantitative, but also
methodological: they make it possible to compare the controlled behavior of the
proxy models with that of the \gls{ai}-based software evaluated in
\textit{black-box} mode. The downstream reporting assets generated under
\codepath{results/figures/chapter_5/} provide the thesis-ready figures,
figure manifests, and summary files derived from these consolidated metrics.


% ─────────────────────────────────────────────────────────────────────────────
\section{Forensic Evaluation Bundle}
\label{app:forensic-evaluation-bundle}
% ─────────────────────────────────────────────────────────────────────────────

The \textit{forensic evaluation bundle} constitutes the operational artifact
used to submit a common set of images to the \gls{ai}-based software evaluated
in \textit{black-box} mode, without revealing the corresponding ground truth in
advance. Its methodological function is to prevent file names, directory
structure, or experimental metadata from influencing the analysis procedure or
making the true class and applied transformation immediately recognizable. The
bundle is generated by the script
\codepath{datasets/scripts/bundle/16_build_forensic_evaluation_bundle.py} and is
saved in the \codepath{datasets/forensic_evaluation_bundle/} directory; its
structure strictly separates the blind inputs intended for the software systems
from the metadata required for the subsequent evaluation through hash-based
mapping.

\begin{lstlisting}[caption={Logical structure of the forensic evaluation bundle.}, label={lst:forensic-evaluation-bundle-structure}]
datasets/forensic_evaluation_bundle/
  blind_tool_input/
    files/
  metadata/
    bundle_manifest.csv
    bundle_hashes_sha256.csv
    bundle_summary.json
  structured_audit_view/
\end{lstlisting}

In the \textit{black-box} evaluation, only the files contained in
\codepath{blind_tool_input/files/} must be imported into the considered
software systems. The \codepath{metadata/} and
\codepath{structured_audit_view/} directories must not be imported, since they
contain information that would reveal true classes, sources, attacks,
transformations, and experimental mapping. This separation reduces the risk of
bias induced by file names, directory structure, or prior knowledge of the
ground truth.

The current bundle includes clean images, \gls{ood} samples, adversarial
perturbations, and anti-forensic transformations. The operational summary of the
bundle documents 11\,500 files in total: 1\,000 clean images, 500 \gls{ood}
samples, 5\,000 adversarial samples, and 5\,000 anti-forensic samples. The
validation checks include uniqueness of the \texttt{bundle\_id}, uniqueness of
the actual hashes, correspondence of hashes when present in the manifest,
semantic cleanliness of the blind paths, and separation of metadata from the
inputs intended for the software systems.


% ─────────────────────────────────────────────────────────────────────────────
\section{Software Environment and Dependencies}
\label{app:software-environment}
% ─────────────────────────────────────────────────────────────────────────────

The pipeline is implemented primarily in Python and uses libraries for tabular
manifest management, image processing, model execution, and production of
experimental outputs. Table~\ref{tab:software-environment} summarizes the main
software components. The exact versions of the dependencies should be maintained
in the repository environment files or in the experimental logs, since
computational reproducibility may be affected by variations in libraries,
numerical backends, \gls{gpu} drivers, and hardware availability.

\begin{table}[H]
\centering
\caption{Main software components used in the consolidated pipeline.}
\label{tab:software-environment}
\begin{tabularx}{\textwidth}{@{}lX@{}}
\toprule
\textbf{Component} & \textbf{Main function} \\
\midrule
\texttt{Python} & Main language used to orchestrate the experimental pipeline. \\
\texttt{pandas} & Management of manifests, metadata, tabular reports, and experimental results. \\
\texttt{NumPy} & Numerical support for arrays, transformations, and matrix operations. \\
\texttt{SciPy} & Support for numerical operations and transformations used in the experimental pipeline. \\
\texttt{Pillow} & Image opening, validation, conversion, and saving. \\
\texttt{PyTorch} & Training, inference, and perturbation generation on neural models. \\
\texttt{torchvision} & \gls{cnn} architectures, preprocessing, and support for proxy model evaluation. \\
\texttt{open\_clip\_torch} & Implementation of the \gls{clip} model used as a proxy model. \\
\texttt{Captum} & Generation of attribution maps through \gls{ig}. \\
\texttt{matplotlib} & Production of plots, curves, experimental figures, and diagnostic visualizations. \\
\bottomrule
\end{tabularx}
\end{table}

The operational-forensic evaluation also includes four selected software systems
observed in \textit{black-box} mode: \gls{magnetaxiomtool}, version
10.1.0.48673, with \gls{magnetai} functionality; \gls{excirefoto2025}, version
4.1.5, used as a stand-alone general-purpose application with \gls{ai}-based
functionality; \gls{cellebriteinseyets}, version 10.9, with exports produced in
the \textit{Cellebrite Physical Analyzer} 10.9.0.3029 environment; and
\gls{griffeye} x64, version 26.2.108, with the \gls{t3kcore} v1.18.0 module.
Unlike Python scripts and proxy models, these systems introduce a component that
is not fully automatable, due to licenses, graphical interfaces, software
versions, operational configurations, and export formats. For this reason, their
evaluation is documented through the blind bundle, operational instructions,
normalized exports, and hash-based mapping.


% ─────────────────────────────────────────────────────────────────────────────
\section{Logs, Reports, and Audit Outputs}
\label{app:logs-reports-audit-outputs}
% ─────────────────────────────────────────────────────────────────────────────

In addition to the main manifests, the repository preserves logs, summaries, and
intermediate reports that are useful for verifying the operations performed.
These artifacts do not all have the same role: some describe the status of the
manual selection process, others summarize split generation, and others document
the composition of the bundle or the evaluation of proxy models. Their overall
function is to make the experimental path directly observable and to support any
subsequent checks.

\begin{lstlisting}[caption={Schematic example of output from the dataset preparation phase.}, label={lst:dataset-preparation-log}]
[INFO] Scanning raw dataset sources...
[INFO] Valid images detected: ...
[INFO] Invalid images discarded: ...
[INFO] SHA256 hashes computed: ...
[INFO] Global duplicates removed: ...
[INFO] Prepared image pool written to datasets/prepared/final_pool/images/
[INFO] Metadata written to datasets/prepared/final_pool/metadata.csv
\end{lstlisting}

The most relevant audit artifacts include
\codepath{manual_selection_log.csv}, \codepath{manual_selection_state.json},
\codepath{manual_selection_summary.json},
\codepath{split_generation_summary.json}, the attack manifests,
\codepath{bundle_summary.json},
\codepath{results/figures/chapter_5/chapter5_figures_manifest.csv},
\codepath{results/figures/chapter_5/chapter5_figures_summary.json}, and
\codepath{results/figures/chapter_5/embedded_metadata_sensitivity_summary.json}.
These files make it possible to reconstruct not only the final output, but also
the status and operational assumptions of the main phases of the pipeline.


% ─────────────────────────────────────────────────────────────────────────────
\section{Reproducibility Boundaries}
\label{app:reproducibility-boundaries}
% ─────────────────────────────────────────────────────────────────────────────

The reproducibility of the pipeline must be interpreted consistently with the
nature of the case study. The deterministically re-executable phases include the
technical preparation of the dataset, manifest generation, split construction,
production of perturbed outputs, and construction of the bundle, within the
limits imposed by software dependencies and input availability. The final
selection of the dataset, instead, includes a deliberate, documented, and
motivated manual component, which must not be concealed as if it were a fully
automatic procedure. The scientific value of this component lies in its
traceability, in the explicit definition of the criteria, and in the freezing of
the final manifest.

Similarly, the evaluation with selected black-box software systems is not
entirely reproducible through command-line execution, since it depends on
proprietary software, specific versions, interface settings, internal
\gls{ai}-based modules, and export procedures. To mitigate this limitation, the
pipeline uses a common blind input, preserves file hashes, separates the ground
truth from the inputs submitted to the software systems, and includes a
normalization phase for the exports. In this way, reproducibility is not
understood as a bit-for-bit replica of every graphical interaction, but as
verifiability of the protocol, inputs, mappings, and final metrics.


% ─────────────────────────────────────────────────────────────────────────────
\section{Repository Availability}
\label{app:repository-availability}
% ─────────────────────────────────────────────────────────────────────────────

The repository associated with the thesis is available at the following address:

\begin{center}
\url{https://github.com/lmolinario/msc-thesis-ai-robustness-in-digital-forensics}
\end{center}

The repository contains the main scripts, official manifests, operational
documentation, and artifacts required to reconstruct the experimental workflow.
Before the final public release, the availability of raw images, generated
outputs, and sensitive materials must be assessed according to the legal,
ethical, and licensing constraints of the original sources. The source code of
the repository is released under the MIT license, as indicated in the
\codepath{LICENSE} file. This license applies to the code and experimental
scripts, but it does not automatically imply the free redistribution of the
thesis text, figures, datasets, raw images, selected software exports, model
checkpoints, or generated artifacts. These materials must be considered
separately, according to their respective legal, ethical, institutional, and
source-dependent licensing constraints.
